% Provenance: P2 abstract/intro, main_zvf.tex framing, advisor brief 2026-07-12,
% consolidation plan v3 claim contract.
\chapter{Introduction}
\label{ch:intro}

\section{The problem: silent signal starvation}
Reinforcement learning from verifiable rewards has become the standard second
stage of language-model training for reasoning tasks. The dominant family of
methods --- GRPO and its variants --- avoids a learned critic by sampling a
\emph{group} of $G$ completions per prompt and centring each completion's
reward against its group mean. This design has a structural blind spot: when
every completion in a group receives the same reward, the centred advantages
are identically zero and the group contributes \emph{no gradient}. Under
binary rewards this happens in exactly two situations --- the model fails
every attempt (too hard) or solves every attempt (mastered). Both ends of
training are therefore starved of signal, and, critically, the second kind of
starvation is invisible in the reward curve: mean reward reads ``success''
precisely while learning has stopped.

This thesis names, measures, calibrates, and budgets that phenomenon through a
single statistic:

\begin{definition}[Zero-Variance Fraction]
At training step $t$, with prompt set $\mathcal{P}_t$ and $G$ completions per
prompt,
\[
\ZVF_t \;=\; \frac{1}{|\mathcal{P}_t|}\sum_{p\in\mathcal{P}_t}
\mathbf{1}\!\left[\operatorname{Var}_{c\sim\mathcal{C}(p)} r(c) = 0\right],
\qquad \GU_t \;=\; 1-\ZVF_t ,
\]
the fraction of groups carrying zero advantage, and its complement, the
fraction of groups carrying usable gradient signal (\emph{gradient
utilisation}).
\end{definition}

\section{Claim contract}
\label{sec:claims}
Following an external adversarial review of this program (Chapter~9), the
thesis is organised around an explicit contract separating what the evidence
supports from what it does not.

\paragraph{Claim 1 (supported; Chapters 4--5).} \ZVF{}, read together with
mean reward, is a near-zero-cost online diagnostic that identifies
zero-advantage regimes. \ZVF{} alone aliases mastery with incapacity; the
reward coordinate disambiguates. Its estimator admits a calibrated confidence
interval (Wilson; empirical coverage 0.95--0.98 in all tested settings), and a
geometric waiting-time model built on it yields a \emph{reliability budget} ---
the $(1{-}\delta)$-quantile of rollouts until the next informative group ---
that matched observed quantiles at ratio 1.00 across six difficulty strata.

\paragraph{Claim 2 (supported; Chapter 6).} At a matched budget of 2{,}560
rollouts per arm on Qwen3-8B/GSM8K (LoRA rank 4), $G{=}2\times160$ steps and
$G{=}16\times20$ steps reach comparable final train reward, but every $G{=}2$
arm terminates inside the all-correct zero-variance wall
($\ZVF\approx0.75$--$1.0$ at reward $\approx1.0$) while $G{=}16$ arms hold
$\ZVF\le0.25$ throughout (seeds 123/456). Group size controls \emph{which end
of training starves}; it is a schedule question, not a constant.

\paragraph{Non-claims (explicit).}
\begin{enumerate}[leftmargin=1.6em,itemsep=1pt]
  \item No predictive or causal power is claimed for \ZVF{} beyond diagnosis;
        in our corpus \ZVF{} rises \emph{after} reward plateaus in every
        collapse we measured.
  \item No data-dependent optimal group size is claimed. Chapter~5 shows the
        signal-per-rollout objective's optimum is the universal
        $G^\star\in\{2,3\}$ for \emph{every} difficulty prior --- an algebraic
        identity discovered during external review of our own theory, reported
        here as an honest negative result.
  \item No adaptive-controller efficacy is claimed. The controller design of
        Chapter~6 is presented with its evidence gaps; the decisive
        pre-registered, compute-matched comparison against static $G{=}16$ and
        naive heuristics is future work.
  \item No generalisation is claimed beyond the studied stack: Qwen3-8B,
        GSM8K-family binary-reward tasks, and a managed LoRA-only training API
        with a closed loss kernel, mostly at 1--3 seeds.
\end{enumerate}

\section{Contributions}
\begin{enumerate}[leftmargin=1.6em,itemsep=1pt]
  \item A measurement study of \ZVF{} across libraries, model scales
        (0.6B--671B logged runs), and training phases (Chapter~4).
  \item Calibrated estimation and a corrected waiting-time reliability budget
        for \ZVF{} (Chapter~5), including two theory corrections found by
        adversarial review and adopted openly.
  \item The matched-budget group-size result (Claim~2) and its consequences
        for group-size scheduling (Chapter~6).
  \item A six-arm uncapped loss-form panel: no length inflation and no
        late-\ZVF{} separation between GRPO and Dr.GRPO at this scale
        (Chapter~7).
  \item A reproducibility contribution grounded in the program's own failures:
        an incident case study, an eight-item minimum-reportable-stack
        standard, a machine-readable registry, and a survival-audit protocol
        (Chapter~9).
  \item A released benchmark harness and artifact trail: run manifests,
        result JSONs, W\&B mirrors, and per-claim provenance
        (Chapter~3, Appendix~B).
\end{enumerate}

\section{Thesis organisation}
Chapter~2 reviews the GRPO family and the reproducibility literature.
Chapter~3 describes the experimental infrastructure. Chapters~4--8 present the
technical core: diagnostic evidence, theory, group size, loss form, and
scaling observations. Chapter~9 presents the reproducibility programme.
Chapter~10 concludes. Appendix~A maps every one of the seventeen working
papers produced during the programme to the chapter that absorbs it;
Appendix~B maps claims to artifacts.
